
Our work bridges three research communities---scaling laws, data curation practices, and data-centric AI---by providing a quantitative framework for data quality that was missing from all three.

\subsubsection{Scaling Laws for Foundation Models}

Neural scaling laws predict model performance from compute budget and architectural choices. Kaplan et al. (2020) showed power-law relationships between validation loss and model size N, dataset size D, and compute C. Hoffmann et al. (2022) refined this with the Chinchilla scaling laws, finding that the compute-optimal ratio is approximately $N^0$.5 tokens per parameter. This led to smaller, better-trained models like Chinchilla (70B parameters, 1.4T tokens) outperforming Gopher (280B parameters, 300B tokens).

However, both formulations assume uniform data quality: the loss function L(N, D, C) treats dataset size D as token count without quality weighting. One trillion deduplicated tokens and one trillion with 40\% duplicates receive identical treatment. While Hoffmann et al. acknowledge that ''data quality matters'' in discussion, they provide no mechanism to incorporate it into the scaling law.

Our contribution: We propose Q(D) as a measurable dimension, demonstrating that quality explains 61\% of information density variance orthogonally to dataset size.

\subsubsection{Data Curation in Foundation Model Pretraining}

Practitioner knowledge about data quality is extensive but heuristic. GPT-3 employed fuzzy deduplication to remove near-duplicate documents. The Pile assembled 22 curated data sources emphasizing domain diversity. LLaMA combined aggressive quality filters with careful corpus mixing ratios. RedPajama reproduced LLaMA's filtering pipeline at scale, documenting heuristics like ''remove documents with >30\% duplicate n-grams.''

Despite this wealth of empirical practice, no prior work has validated which quality dimensions matter most via controlled experiments. Is deduplication more important than diversity? Does perplexity filtering provide value beyond removing non-English text? Practitioners lack quantitative evidence to prioritize curation investments.

Our contribution: We isolate and measure four quality dimensions on controlled C4 subsets, demonstrating that deduplication is the strongest predictor (r = 0.72) while diversity (r = 0.65), perplexity (r = 0.58), and efficiency (r = 0.53) all contribute non-redundantly.

\subsubsection{Data-Centric AI and Quality Metrics}

The data-centric AI movement argues that improving data often yields larger performance gains than architectural innovations. Northcutt et al. (2021) showed that label noise degrades model performance and proposed automated cleaning methods. Sorscher et al. (2022) demonstrated that data pruning can match full-dataset performance with 50\% fewer training samples. Abbas et al. (2023) showed 40\% data reduction via perplexity-based pruning with minimal loss degradation.

However, this work remains qualitative about ''quality'' itself. Data pruning methods define quality as ''samples the current model finds easy/hard'' (perplexity-based), which is task-specific and requires a trained model---not suitable for pretraining from scratch.

Our contribution: We define quality via information density---an intrinsic, task-agnostic property measurable before training begins. Our Q(D) metric combines deduplication, diversity, perplexity, and efficiency into a composite score validated against information-theoretic ground truth. Unlike pruning methods that require a trained model, Q(D) can be computed on raw text.

\subsubsection{Information Theory and Learning Efficiency}

Our work connects to information-theoretic perspectives on learning efficiency. Shannon entropy quantifies information content in data. Rate-distortion theory predicts that compression improves signal transmission efficiency. Recent work has applied information theory to dataset quality. Ash et al. (2020) used gradient-based scores to select informative training examples. However, these methods measure ''informativeness to a specific model'' rather than intrinsic data quality.

Our contribution: We measure information density via entropy and redundancy before any model training, providing a model-agnostic quality metric.

\subsubsection{Positioning Summary}

Our work occupies the intersection of three research areas: we provide the quantitative quality metric that scaling laws lack, the controlled experimental validation that practitioner heuristics lack, and the task-agnostic measurement framework that data-centric AI lacks. We are the first to validate a multi-component Q(D) metric against information density through controlled experiments suitable for scaling law integration.
